\chapter{La planification CORA en 3D}
\label{chap:planification}

\begin{chapeau}
C'est le cœur du projet. À partir de l'os et de ses repères, la planification construit les axes des deux segments,
trouve le CORA et l'angle vrai, place la charnière et les traits de coupe, calcule la correction (avec la dérotation),
place les broches, choisit la voie d'abord et le type de coin, puis écrit la prescription. Tout est calculé dans le
fichier JavaScript \fichier{code_source/cora/app/cora_simple.js} (fonction \code{planifier()}), qui ne fait aucun
dessin : il peut être testé seul, sous Node.
\end{chapeau}

\section{Les axes des deux segments}
\label{sec:axes}

L'os déformé est vu comme deux segments droits. Il faut l'axe de chacun. Paley le construit à partir de la ligne
articulaire de l'extrémité et de l'angle articulaire \emph{normal} \cite{paley2002}. Le logiciel offre deux
références.

\subsection{Référence « os sain »}

Si l'on dispose d'un os sain du même animal (le membre controlatéral), ou de l'os avant une déformité de test, on
mémorise la position de son axe par rapport à ses articulations, puis on la reporte sur l'os déformé. Pour une
articulation, on construit un repère local avec ses deux lignes articulaires :
\begin{equation}
  \vb{e}_1 = \frac{\vb{J}_F}{\norm{\vb{J}_F}}, \qquad
  \vb{e}_3 = \frac{\vb{J}_S - (\vb{J}_S\cdot\vb{e}_1)\,\vb{e}_1}{\norm{\ldots}}, \qquad
  \vb{e}_2 = \vb{e}_3 \times \vb{e}_1,
\end{equation}
où $\vb{J}_F$ est la ligne articulaire de face (de médial à latéral) et $\vb{J}_S$ celle de profil (de crânial à
caudal). L'axe sain s'écrit dans ce repère local ; sur l'os déformé, on reconstruit le même repère local avec ses
propres lignes articulaires et on y replace l'axe. C'est l'équivalent 3D de « ligne articulaire + angle normal ».

\subsection{Référence « valeurs normales » (méthode de Paley)}

Sans os sain, on construit l'axe dans chaque vue. Dans la vue de face, on tourne la direction de la ligne articulaire
de l'angle normal, dans un sens puis dans l'autre, et on garde la direction qui va vers distal et qui est la plus
proche de la verticale. On fait de même de profil. On obtient deux directions 2D, $\vb{d}_F = (d_{F,x}, d_{F,y})$ de
face et $\vb{d}_S = (d_{S,z}, d_{S,y})$ de profil, qu'on combine en une direction 3D dont elles sont les projections :
\begin{equation}
  \vb{d} = \frac{\bigl(d_{F,x}/\abs{d_{F,y}},\ -1,\ d_{S,z}/\abs{d_{S,y}}\bigr)}{\norm{\ldots}}.
  \label{eq:dir3d}
\end{equation}
Au tibia, les axes passent par l'éminence (segment proximal) et par le centre distal (segment distal). Au fémur, l'axe
proximal est l'axe anatomique mesuré ; l'axe distal a la direction donnée par les lignes articulaires distales, et sa
position est fixée de face par l'os fictif (il coupe la ligne des condyles au même endroit que l'axe sain), et de
profil par la règle « une seule angulation » : le CORA de profil est au même niveau que le CORA de face.

\begin{pourquoi}
\textbf{Pourquoi cette règle de profil au fémur ?} Le fémur du chien est arqué vers l'avant : chez le lévrier, les
condyles sont 33~mm en arrière de l'axe proximal. Un gabarit de profil tiré d'un autre os serait donc faux d'un chien à
l'autre. Pour une déformité à un seul apex, les CORA de face et de profil sont au même niveau : c'est une hypothèse
plus sûre.
\end{pourquoi}

\section{Le CORA en 3D}
\label{sec:cora3d}

\subsection{Deux droites qui ne se coupent pas}

En 3D, deux axes mesurés ne se coupent presque jamais exactement : ils passent l'un à côté de l'autre. On cherche le
point où ils sont le plus proches. Soient l'axe proximal $\vb{p} + s\,\vb{u}$ et l'axe distal $\vb{q} + t\,\vb{v}$
(\cref{fig:perp-commune}). On minimise
\begin{equation}
  f(s, t) = \norm{\vb{w}_0 + s\,\vb{u} - t\,\vb{v}}^2, \qquad \vb{w}_0 = \vb{p} - \vb{q}.
\end{equation}
En annulant les dérivées par rapport à $s$ et à $t$, on obtient un système $2 \times 2$. Avec $a = \vb{u}\cdot\vb{u}$,
$b = \vb{u}\cdot\vb{v}$, $c = \vb{v}\cdot\vb{v}$, $d = \vb{u}\cdot\vb{w}_0$ et $e = \vb{v}\cdot\vb{w}_0$ :
\begin{equation}
  \tcbhighmath{s^* = \frac{b\,e - c\,d}{a\,c - b^2}, \qquad t^* = \frac{a\,e - b\,d}{a\,c - b^2}.}
  \label{eq:perp}
\end{equation}
Les points $P_p = \vb{p} + s^*\vb{u}$ et $P_d = \vb{q} + t^*\vb{v}$ sont les pieds de la \emph{perpendiculaire
commune}. Le CORA est leur milieu, $C = (P_p + P_d)/2$, et $\delta = \norm{P_p - P_d}$ dit de combien les axes se
ratent. Si $a\,c - b^2 \approx 0$, les axes sont parallèles : il n'y a pas d'angulation.

\fig{perpendiculaire_commune.pdf}{Le CORA en 3D : le milieu de la perpendiculaire commune aux deux
axes.}{fig:perp-commune}

\subsection{Quand il n'y a pas de CORA}

Si l'angle entre les axes est inférieur à 0,5\degre, ou si le CORA tombe hors de l'os (au-delà de 4 à 96~\% de sa
longueur), il n'y a pas d'ostéotomie au CORA possible. C'est le cas des os sains du lévrier : le logiciel l'indique
(« pas de correction indiquée ») au lieu de proposer une coupe absurde.

\section{L'angle vrai et le plan de la déformité}
\label{sec:angle-vrai}

L'angle vrai est simplement l'angle entre les deux directions :
\begin{equation}
  \theta = \arccos(\vb{u}\cdot\vb{v}).
\end{equation}
Les composantes $F$ et $S$ sont les angles orientés entre les projections de $\vb{u}$ et $\vb{v}$ dans chaque vue. De
face, $F < 0$ indique un varus (le bas part vers médial) ; de profil, $S > 0$ un procurvatum.

\textbf{Démonstration de $\tan^2\theta = \tan^2 F + \tan^2 S$.} Prenons l'axe proximal vertical,
$\vb{u} = (0, -1, 0)$. Une direction distale dont les angles projetés valent $F$ de face et $S$ de profil a pour
vecteur directeur $\vb{v} \propto (\pm\tan F,\ -1,\ \pm\tan S)$ : de face, on ne voit que $(\pm\tan F, -1)$, qui fait
bien l'angle $F$ avec la verticale ; de profil, $(\pm\tan S, -1)$. Alors
\begin{equation}
  \tcbhighmath{\cos\theta = \frac{\vb{u}\cdot\vb{v}}{\norm{\vb{v}}} = \frac{1}{\sqrt{1 + \tan^2 F + \tan^2 S}}
  \quad\Longrightarrow\quad
  \tan^2\theta = \frac{1}{\cos^2\theta} - 1 = \tan^2 F + \tan^2 S.}
\end{equation}
Le logiciel n'a pas besoin de cette formule, puisqu'il calcule $\theta$ directement en 3D ; elle sert de contrôle et
d'explication (\cref{fig:angle-vrai}). L'orientation du plan de la déformité par rapport au plan frontal est
$\varphi = \arctan\bigl(\abs{\tan S}/\abs{\tan F}\bigr)$.

\fig{angle_vrai.pdf}{À gauche, les projections $F$ et $S$ et l'angle vrai $\theta$. À droite, l'angle vrai en
fonction de $F$ pour plusieurs valeurs de $S$ : il est toujours inférieur à la somme $F + S$ (en tirets).}{fig:angle-vrai}

\section{La bissectrice, l'ACA et la rotation}
\label{sec:bissectrice}

Les deux axes définissent le plan de la déformité. Trois directions servent ensuite (\cref{fig:bissectrice}a) :
\begin{equation}
  \vb{n}_{\text{ACA}} = \frac{\vb{u}\times\vb{v}}{\norm{\vb{u}\times\vb{v}}}, \qquad
  \vb{b} = \frac{\vb{u} + \vb{v}}{\norm{\vb{u} + \vb{v}}}, \qquad
  \vb{c} = \frac{\vb{u} - \vb{v}}{\norm{\vb{u} - \vb{v}}}.
\end{equation}
$\vb{n}_{\text{ACA}}$ est perpendiculaire au plan de la déformité : c'est la direction de l'axe de correction.
$\vb{b}$ est la normale du \emph{plan bissecteur}, qui partage l'angle entre les segments. $\vb{c}$ est la
\emph{bissectrice transverse}, dans le plan de la déformité, dirigée vers le côté convexe. Le logiciel coupe l'os par
le plan bissecteur passant par le CORA : il en tire la largeur de l'os le long de $\vb{c}$ (distances aux corticales
convexe et concave) et sa profondeur le long de l'ACA.

\subsection{La formule de Rodrigues}

La correction fait tourner le segment distal d'un angle $\theta$ autour de l'axe unitaire
$\vb{k} = (\vb{v}\times\vb{u})/\norm{\vb{v}\times\vb{u}}$, qui amène $\vb{v}$ sur $\vb{u}$. La matrice de rotation
d'angle $\alpha$ autour de $\vb{k}$ est donnée par la formule de Rodrigues \cite{rodrigues1840} :
\begin{equation}
  \tcbhighmath{\mt{R}(\vb{k}, \alpha) = \mt{I} + \sin\alpha\,[\vb{k}]_\times + (1 - \cos\alpha)\,[\vb{k}]_\times^2,
  \qquad
  [\vb{k}]_\times = \begin{pmatrix} 0 & -k_z & k_y \\ k_z & 0 & -k_x \\ -k_y & k_x & 0 \end{pmatrix}.}
  \label{eq:rodrigues}
\end{equation}
Une rotation autour d'une droite qui passe par le point $Q$ envoie $\vb{x}$ sur
$\mt{R}(\vb{x} - Q) + Q$ : c'est une transformation rigide $\mt{T} = (\mt{R}, \vb{t})$ avec $\vb{t} = Q - \mt{R}Q$.
Deux transformations se composent par $\mt{T}_A \circ \mt{T}_B = (\mt{R}_A\mt{R}_B,\ \mt{R}_A\vb{t}_B + \vb{t}_A)$, et
l'inverse est $(\mt{R}^{\mathsf T}, -\mt{R}^{\mathsf T}\vb{t})$.

\section{L'ostéotomie : charnière, traits et coin}
\label{sec:traits}

\subsection{La charnière}

La charnière $H$ est placée sur la bissectrice transverse, au niveau du CORA, à une profondeur $h$ (2~mm par défaut)
à l'intérieur de la corticale opposée au coin :
\begin{equation}
  H = P_c + \sigma\,(\ell - h)\,\vb{c},
  \qquad
  \begin{cases}
    \sigma = -1,\ \ell = \text{distance à la corticale concave} & \text{coin de fermeture,} \\
    \sigma = +1,\ \ell = \text{distance à la corticale convexe} & \text{coin d'ouverture,}
  \end{cases}
\end{equation}
où $P_c$ est le centre de la section. La lame entre de l'autre côté : par la corticale convexe pour une fermeture, par
la corticale concave pour une ouverture. On note $D = \text{largeur} - h$ la distance de la charnière à la corticale
d'entrée.

\begin{pourquoi}
\textbf{Pourquoi 2~mm ?} À cet endroit, la charnière est l'ACA de Paley : aucune translation. Et il reste un pont d'os
cortical assez fin pour plier sans casser (au plus 3~mm), qui garde les fragments en contact. Au-delà de 3~mm, le pont
ne plie plus : l'ostéotomie devient complète et le coin devient \emph{mixte} (une partie se ferme, l'autre s'ouvre).
\end{pourquoi}

\subsection{Les traits de coupe}

\begin{itemize}
  \item \textbf{Coin de fermeture} : deux plans passant par $H$, de normales $\vb{u}$ et $\vb{v}$ (l'un est
  perpendiculaire à l'axe proximal, l'autre à l'axe distal). Le coin d'os entre les deux est retiré.
  \item \textbf{Coin d'ouverture} : un seul plan passant par $H$, de normale $\vb{b}$ (le plan bissecteur).
\end{itemize}

\fig{bissectrice_coins.pdf}{(a) Le plan de la déformité : axes, CORA, bissectrice transverse. (b) Le coin de fermeture
et sa base. (c) Le coin d'ouverture et l'écart qu'il crée.}{fig:bissectrice}

\subsection{La taille du coin}

\textbf{Fermeture.} Prenons un point à la distance $D$ de $H$ le long de $\vb{c}$, et mesurons l'écart entre les deux
plans le long de $\vb{b}$. Un point $H + D\,\vb{c} + w\,\vb{b}$ est sur le plan de normale $\vb{u}$ si
$D\,(\vb{c}\cdot\vb{u}) + w\,(\vb{b}\cdot\vb{u}) = 0$. Or $\vb{c}\cdot\vb{u} = \sin(\theta/2)$ et
$\vb{b}\cdot\vb{u} = \cos(\theta/2)$, d'où $w = -D\tan(\theta/2)$ ; de même $w = +D\tan(\theta/2)$ pour le plan de
normale $\vb{v}$. La base du coin, à la corticale d'entrée, vaut donc
\begin{equation}
  \tcbhighmath{\text{base}_{\text{fermeture}} = 2\,D\,\tan(\theta/2).}
\end{equation}

\textbf{Ouverture.} Un point de la face coupée, à la distance $D$ de la charnière, décrit un arc de cercle d'angle
$\theta$ : la corde de cet arc est l'ouverture
\begin{equation}
  \tcbhighmath{\text{ouverture} = 2\,D\,\sin(\theta/2).}
\end{equation}

\begin{exemple}
Tibia de 16~mm de large, charnière à 2~mm, $\theta = 20\degre$ : $D \approx 14$~mm, base du coin de fermeture
$2 \times 14 \times \tan 10\degre \approx 4{,}9$~mm, ouverture $2 \times 14 \times \sin 10\degre \approx 4{,}9$~mm. Pour un
petit angle, les deux sont proches. Il faut au moins 1,5~mm de matière entre les deux fentes du guide : un coin de
fermeture trop fin est déconseillé (\cref{sec:choix-coin}).
\end{exemple}

\subsection{Rester loin du cartilage}

Un trait de coupe doit rester à au moins 5~mm du cartilage articulaire, et un coin de fermeture ne doit pas en
emporter. Si le CORA est trop près de l'articulation (par exemple dans la trochlée du fémur), le logiciel déplace la
coupe vers la diaphyse, millimètre par millimètre, jusqu'à 35~\% de la longueur de l'os : c'est la règle 2.

\section{La règle 2 : couper ailleurs qu'au CORA}
\label{sec:regle2}

Quand la coupe est déplacée de $z$~mm le long de l'os, la charnière $H$ n'est plus sur la bissectrice du CORA. Paley
impose alors que l'\textbf{ACA reste sur la bissectrice} : sinon les axes ne seraient pas réalignés (règle 3).

\textbf{Pourquoi tout point de la bissectrice convient.} Soit $Q = C + \lambda\,\vb{c}$ un point de la bissectrice et
$\mt{R}$ la rotation d'angle $\theta$ qui amène $\vb{v}$ sur $\vb{u}$. Elle envoie la direction distale sur
$\vb{u}$ ; il reste à vérifier que l'axe distal tombe sur l'axe proximal. Le point $C$, commun aux deux axes, va en
$C' = C + \lambda\,(\vb{c} - \mt{R}\vb{c})$. Dans le plan de la déformité, $\vb{v}$ est $\vb{u}$ tourné de $-\theta$
et $\mt{R}\vb{u}$ est $\vb{u}$ tourné de $+\theta$, donc $\vb{v} + \mt{R}\vb{u} = 2\cos\theta\,\vb{u}$. Comme
$\vb{c} \propto \vb{u} - \vb{v}$ et $\mt{R}\vb{c} \propto \mt{R}\vb{u} - \vb{u}$ :
\begin{equation}
  \vb{c} - \mt{R}\vb{c} \propto 2\vb{u} - (\vb{v} + \mt{R}\vb{u}) = 2(1 - \cos\theta)\,\vb{u}.
\end{equation}
$C'$ est donc sur l'axe proximal, et l'axe distal corrigé est confondu avec lui, quel que soit $\lambda$.

\textbf{Le bon point.} Parmi ces points, le logiciel prend celui qui garde les faces de l'ostéotomie en contact : la
charnière doit rester sur le plan de la face proximale (de normale $\vb{n}_F$ : $\vb{u}$ en fermeture, $\vb{b}$ en
ouverture). La condition $(\mt{R}(H - Q) + Q - H)\cdot\vb{n}_F = 0$ donne
\begin{equation}
  \lambda = \frac{\bigl((\mt{R} - \mt{I})(H - C)\bigr)\cdot\vb{n}_F}{\bigl((\mt{R} - \mt{I})\,\vb{c}\bigr)\cdot\vb{n}_F}.
\end{equation}
Le fragment distal glisse alors le long de la coupe : c'est la \textbf{translation à l'ostéotomie}, mesurée par
$\norm{\mt{T}(H) - H}$ (\cref{fig:regle2}). Au-delà de 50~\% de la largeur de l'os, le contact osseux devient
insuffisant : le logiciel le signale comme bloquant.

\fig{regle2.pdf}{Règle 1 et règle 2 (coin d'ouverture). En coupant 22~mm sous le CORA, les axes sont toujours
réalignés, mais le fragment distal (en bleu) glisse le long de la coupe.}{fig:regle2}

\section{La dérotation}
\label{sec:derotation}

Si la torsion doit être corrigée, le fragment distal tourne en plus autour de son axe, d'un angle $\alpha$. La
\textbf{cible} dépend de la référence : la torsion de l'os sain, ou, avec les valeurs normales, la valeur moyenne
publiée seulement si la torsion sort des seuils. La décision se prend sur la torsion \emph{après} la seule correction
angulaire : une angulation pure incline un peu les axes transverses et crée quelques dixièmes de degré de « fausse
torsion », que la correction annule.

La dérotation tourne autour d'un axe que la correction garde aligné :
\begin{itemize}
  \item \textbf{fermeture} : autour de l'axe distal, \emph{avant} la fermeture, $\mt{T} = \mt{T}_{\text{angle}} \circ
  \mt{R}_\alpha$ ; la face distale, perpendiculaire à cet axe, tourne dans son plan ;
  \item \textbf{ouverture} : autour de l'axe réaligné (l'axe proximal), \emph{après} l'ouverture,
  $\mt{T} = \mt{R}_\alpha \circ \mt{T}_{\text{angle}}$.
\end{itemize}

L'angle $\alpha$ est cherché par la \textbf{méthode de la sécante} \cite{press2007}, pour que la torsion mesurée après
la correction, $\tau(\alpha)$, égale la cible $\tau^*$. Avec $f(\alpha) = \tau(\alpha) - \tau^*$ :
\begin{equation}
  \alpha_{n+1} = \alpha_n - f(\alpha_n)\,\frac{\alpha_n - \alpha_{n-1}}{f(\alpha_n) - f(\alpha_{n-1})}.
\end{equation}
Comme $\tau$ dépend presque linéairement de $\alpha$, deux ou trois itérations suffisent (au plus six, arrêt à
0,001\degre). On ne dérote pas en dessous de 1\degre. Une dérotation impose une ostéotomie complète : une charnière
corticale ne peut pas tourner sur elle-même.

\section{Les broches}
\label{sec:broches}

Les broches relient la planification au geste : elles fixent le guide, puis servent de repère pour la correction et
de support pour la réduction. Leur direction et leur position sont calculées ici.

\subsection{L'astuce de l'inverse}

Les broches ont quatre rôles : fixer le guide de coupe, porter la correction, tenir la réduction et bloquer la
rotation. L'astuce est de forer les broches distales avec l'\textbf{inverse de la correction}. Si $\vb{d}$ est la
direction des broches proximales et $\mt{R}_T$ la rotation de la correction complète (angle et dérotation), les
broches distales sont forées selon
\begin{equation}
  \vb{d}_{\text{dist}} = \mt{R}_T^{-1}\,\vb{d}
  \quad\Longrightarrow\quad
  \text{après correction : } \mt{R}_T\,\vb{d}_{\text{dist}} = \vb{d}.
\end{equation}
Après la correction, \textbf{toutes les broches sont parallèles}, et seulement si la correction est la bonne. Le
chirurgien a un repère visuel immédiat, et le guide de réduction (\cref{sec:reduction}) peut s'y enfiler.

\fig{broches.pdf}{Les broches distales, forées avec l'inverse de la correction, deviennent parallèles aux broches
proximales quand la correction est exacte.}{fig:broches}

\subsection{Où placer les broches}

Le logiciel cherche une position pour chaque broche dans un grand nombre de possibilités, et garde la première qui
respecte toutes les règles :
\begin{itemize}
  \item \textbf{position le long de l'os} : à une distance de la coupe de 1,2 fois la largeur locale de l'os si la
  place le permet, sinon plus près, jamais contre la tour de fente ;
  \item \textbf{une broche par fragment} (une proximale, une distale), sur la ligne centrale de l'os, ou décalée de 2 à
  4~mm si un obstacle l'impose ;
  \item \textbf{direction} : vers l'os depuis la face d'appui du guide, tournée autour de l'os d'un angle $\beta$ de
  0 à $\pm 45\degre$ (par pas de 5\degre), puis inclinée le long de l'os de $\gamma = \pm 10$ ou $\pm 20\degre$ si le
  retrait du guide l'exige ;
  \item \textbf{hauteur des manchons} : 15~mm, raccourcis à 10 puis 8~mm si un manchon heurte la tour de fente.
\end{itemize}

\begin{table}[htbp]
\centering\small
\caption{Les règles qu'une broche doit respecter.}
\label{tab:regles-broches}
\begin{tabularx}{\linewidth}{L{4.4cm}Y}
\toprule
\titrecol{Règle} & \titrecol{Valeur} \\
\midrule
traverser l'os (bicorticale) & au moins 5~mm d'os le long de la broche \\
rester loin des plans de coupe & 2~mm de pont osseux, du bon côté ; pointe hors du trajet de la lame de l'autre trait \\
rester loin du cartilage & 3~mm de l'axe de la broche ; manchon jamais posé sur le cartilage \\
éviter les os voisins & 2~mm de l'axe, jusqu'à 8~mm au-delà de la corticale opposée (un forage un peu trop long n'embroche pas la fibula) \\
entrer franchement dans la corticale & moins de 50\degre{} entre la broche et la normale à l'os (sinon la mèche glisse) \\
laisser sortir le guide & une des deux moitiés doit s'éloigner de l'autre en sortant (dérive inférieure à 0,03, environ 2\degre) \\
\bottomrule
\end{tabularx}
\end{table}

Deux autres configurations existent : \textbf{1 + 1} (une broche par fragment ; chaque fragment peut alors tourner
autour de sa broche) et \textbf{1 + 2} (une proximale, deux distales côte à côte, utile pour un fragment distal court).

\begin{danslecode}
\code{placerBroches()}, \code{evaluerBroche()} et \code{manchonLibre()}. La recherche évalue des milliers de
positions ; chaque contrôle parcourait tous les sommets de l'os, du cartilage et des voisins (environ 40~s sur un fémur).
Les points sont maintenant rangés une fois dans une grille de cellules de 1,5~mm, et une requête n'examine que les
cellules utiles, avec le même test point par point : mêmes résultats, beaucoup plus vite (\code{grille()},
\code{parcourirLeLong()}).
\end{danslecode}

\section{Où poser le guide : la voie d'abord}
\label{sec:abord}

La lame entre toujours du côté \textbf{opposé à la charnière}, par la \emph{face du coin} : convexe en fermeture,
concave en ouverture. Le guide se pose sur cette face. Mais elle doit être accessible par une voie d'abord habituelle
(\cref{tab:abords}).

\begin{table}[htbp]
\centering\small
\caption{Les secteurs où le guide peut être posé (choix raisonnés du projet, à valider avec un chirurgien).}
\label{tab:abords}
\begin{tabularx}{\linewidth}{lL{3.2cm}Y}
\toprule
\titrecol{Os} & \titrecol{Abord (demi-ouverture)} & \titrecol{Structures à protéger} \\
\midrule
fémur & latéral (± 60\degre) & nerf sciatique en arrière du fémur ; vaste latéral et biceps fémoral \\
 & médial (± 45\degre) & artère et veine fémorales, nerf saphène ; artère génienne descendante \\
tibia & médial (± 60\degre) & veine et artère saphènes, nerf saphène ; derrière l'os, artères poplitée et tibiale crâniale \\
 & latéral (± 30\degre) & muscles tibial crânial et extenseurs ; artère tibiale crâniale et nerf fibulaire \\
\bottomrule
\end{tabularx}
\end{table}

Le logiciel cherche, dans chaque secteur, la direction la plus proche de la face du coin, par pas de 5\degre, telle
qu'aucun os voisin ne soit sur le trajet de la lame. Trois cas :
\begin{itemize}
  \item la face du coin est accessible : le guide s'y pose et la charnière peut être conservée ;
  \item une face accessible est à moins de 60\degre{} de la face du coin : le guide s'y pose, mais l'ostéotomie
  devient complète ;
  \item sinon, la lame entrerait du côté de la charnière : ce coin est déclaré \textbf{non réalisable par un abord
  usuel}, et l'autre type de coin est recommandé. Si en plus un os voisin (en général la fibula) est sur le trajet de
  la lame, \textbf{aucun guide n'est construit} : il occuperait l'espace entre les deux os, où l'on ne peut ni le
  glisser, ni le fixer, ni couper.
\end{itemize}

\section{La stratégie chirurgicale}
\label{sec:strategie}

\subsection{Les contraintes au niveau de la coupe}

Pour chaque coin, le logiciel liste ce qui pourrait empêcher de couper à ce niveau, par gravité (bloquant, attention,
information) :
\begin{itemize}
  \item cartilage et capsule articulaire : au moins 5~mm du trait ;
  \item tubérosité tibiale (insertion du ligament de la rotule) : au moins 8~mm ;
  \item épicondyles du fémur (origine des ligaments collatéraux) : au moins 6~mm ;
  \item physes d'un chien en croissance (zones approchées en fraction de la longueur) ;
  \item place pour la fixation : au moins 25~mm d'os de chaque côté (plaque à trois vis et broches) ;
  \item consolidation : rapide en métaphyse, plus lente en diaphyse et au tiers distal du tibia ;
  \item translation de la règle 2 (attention au-delà de 25~\% de la largeur, bloquant au-delà de 50~\%) ;
  \item structures à protéger pendant l'abord (\cref{tab:abords}).
\end{itemize}

\subsection{Ouverture ou fermeture ?}
\label{sec:choix-coin}

Le logiciel compare les deux coins critère par critère, dans un ordre fixé. Le premier critère qui les départage
décide ; tous restent affichés à l'écran (\cref{fig:choix-coin}) :
\begin{enumerate}
  \item \textbf{accès} : la lame entre par un abord usuel, du côté opposé à la charnière ;
  \item \textbf{charnière conservée} (stabilité, pas de translation) ;
  \item \textbf{règle 1} plutôt que règle 2 ;
  \item un coin de fermeture \textbf{assez épais} pour deux fentes (au moins 1,5~mm de matière entre elles) ;
  \item une \textbf{ouverture raisonnable} (au plus 40~\% de la largeur de l'os ; au-delà, greffe et consolidation
  longue) ;
  \item la \textbf{longueur} : si le membre opéré est plus court (ou plus long) que l'autre d'au moins 3~mm, le coin qui
  compense le mieux (la fermeture raccourcit, l'ouverture allonge) ;
  \item à défaut, la \textbf{fermeture} : surfaces d'os en contact, stabilité immédiate, pas de greffe.
\end{enumerate}

\fig[0.92]{choix_du_coin.pdf}{Comment le logiciel choisit entre le coin de fermeture et le coin
d'ouverture.}{fig:choix-coin}

\begin{exemple}
Sur 606 cas simulés (\cref{sec:campagne-strategie}), les choix suivent la pratique publiée : au tibia, un valgus se
corrige par fermeture et un varus par ouverture (le critère décisif est l'accès : la face médiale du tibia) ; au fémur,
la fermeture domine (critère de stabilité) \cite{petazzoni2012,brower2017}.
\end{exemple}

\section{La prescription}

La planification se termine par un fichier \fichier{prescription.json} (format v2.2), qui contient tout ce que l'étape
7 doit savoir : la déformité, la stratégie, la charnière, les plans de coupe et leur profondeur, la correction (matrice
$4 \times 4$ et dérotation), les broches et leur position après correction, la face d'appui, les contraintes et les
alertes. Il est décrit en détail à l'\cref{ann:prescription}. Pour un os segmenté, il contient aussi la matrice qui
ramène le repère de l'os à celui du scanner : le guide peut ainsi être superposé au STL de l'os dans 3D Slicer.

\section{Les déformités simulées}
\label{sec:deformite-simulee}

Pour tester la méthode, on impose à un os sain une déformité connue (fonction \code{makeDeformity()}). L'apex $A$ est
placé à une fraction choisie de la longueur. La direction distale est
\begin{equation}
  \vb{v}_0 = \frac{(-s_F\tan F,\ -1,\ s_S\tan S)}{\norm{\ldots}},
\end{equation}
avec $s_F = +1$ en varus, $-1$ en valgus, et $s_S = +1$ en procurvatum, $-1$ en recurvatum : ainsi, les angles projetés
valent \emph{exactement} $F$ et $S$. Le segment distal tourne de $\theta = \arccos(\vb{u}_0\cdot\vb{v}_0)$ autour de
$\vb{n} = \vb{u}_0\times\vb{v}_0$ passant par l'apex, après une éventuelle torsion $\tau$ autour de l'axe sain. Pour
éviter une cassure, la rotation est appliquée progressivement dans une zone de transition de largeur $w$ autour de
l'apex, avec le poids
\begin{equation}
  \omega(t) = t^2\,(3 - 2t), \qquad t = \clamp\!\left(\frac{y_A + w/2 - y}{w},\ 0,\ 1\right),
\end{equation}
où $y$ est la hauteur du point. Les 17 \textbf{préréglages cliniques} de l'application (valgus proximal du tibia, pes
varus du teckel, varus du fémur distal de la luxation de la rotule\ldots) sont des déformités de ce type, avec leur
contexte et leur traitement publié.

\section{Les 13 étapes à l'écran}

Dans l'application, la planification est découpée en 13 étapes. Chacune montre l'os en 3D, avec un texte qui
explique ce qui se passe et ses sources (\cref{tab:13etapes}, \cref{fig:capture-angle-vrai}).

\begin{table}[htbp]
\centering\small
\caption{Les 13 étapes de l'onglet \ui{Planification CORA}.}
\label{tab:13etapes}
\begin{tabularx}{\linewidth}{rL{4cm}Y}
\toprule
& \titrecol{Étape} & \titrecol{Ce qu'elle montre} \\
\midrule
1 & Le cas & l'os, ses repères, son côté, le dépistage des angles (± 2 écarts-types) \\
2 & Face : lignes articulaires & les lignes d'orientation articulaire de face \\
3 & Face : angles articulaires & l'axe et les angles de face, comparés aux normes \\
4 & Face : axes et CORA & les axes des segments et le CORA de face \\
5 & Profil : lignes et angles & les lignes articulaires et les angles de profil \\
6 & Profil : axes et CORA & les axes et le CORA de profil \\
7 & Angle vrai (3D) & l'angle vrai et le plan oblique de la déformité \\
8 & Bissectrice et ACA & le plan bissecteur et l'axe de correction \\
9 & Angle et traits de coupe & la charnière, les plans de coupe, le coin \\
10 & Correction & la correction animée, avec la dérotation et les broches \\
11 & Contrôle & les angles après correction, dans les deux plans \\
12 & Guides de coupe & la prescription, la construction des guides, les contrôles, les fichiers \\
13 & Et la torsion ? & la torsion mesurée, et pourquoi la radiographie ne suffit pas \\
\bottomrule
\end{tabularx}
\end{table}

\fig{capture_angle_vrai.png}{L'étape 7 de l'application : un tibia avec 15\degre{} de varus et 10\degre{} de
procurvatum. La déformité est dans un plan oblique (en rose) et l'angle vrai vaut 17,8\degre.}{fig:capture-angle-vrai}

\begin{danslecode}
\fichier{code_source/cora/app/cora_simple.js} : \code{reference()} et \code{axesNormales()} (les deux références),
\code{coraPerpCommune()} (\cref{eq:perp}), \code{planifier()} (tout le calcul), \code{M3.axisAngle()} (Rodrigues),
\code{RT.aboutAxis()} et \code{RT.compose()}, \code{choisirAbord()}, \code{contraintesCoupe()},
\code{recommanderCoin()}, \code{placerBroches()}, \code{mesuresTorsion()}, \code{makeDeformity()},
\code{prescription()}. L'interface des 13 étapes est dans \fichier{code_source/cora/app/modules/40_etapes.js}.
\end{danslecode}
